\chapter{Mathematical Foundations}\label{ch:foundations}

This chapter establishes the notation and background. Section~\ref{sec:nn_notation_pwl} introduces feedforward networks and their \ac{relu}--\ac{cpwl} equivalence. Section~\ref{sec:convex_monotone_foundations} defines monotone networks, \acp{icnn}, and Maxout units, the building blocks of Chapter~\ref{ch:decompositions}. Section~\ref{sec:game_theory_mdp_foundations} presents \acp{mdp}, zero-sum turn-based games, and the full ReLU Net Game of Gaubert and Vlassopoulos, distinguishing the framework inherited by Chapter~\ref{ch:game_theory} from its extensions. The absolute-curvature construction for smooth activations accompanies the network decomposition theorem in Section~\ref{sec:dc_splitting_smooth}.

\section{Neural Network Notation and PWL Mappings}\label{sec:nn_notation_pwl}

The standard feedforward architecture serves as the reference model throughout; convolutional and attention-based architectures are treated in Section~\ref{sec:cnn_decompositions} and Chapter~\ref{ch:transformers} as structured special cases and extensions.

\begin{definition}[Feedforward \ac{relu} network]\label{def:ffn_relu}
Let $L \in \NN$ and let $d^{(0)}, d^{(1)}, \ldots, d^{(L)} \in \NN$ denote the \emph{layer widths}, with input dimension $d^{(0)}$ and output dimension $d^{(L)}$. A \emph{feedforward \ac{relu} network} of depth $L$ is given by weight matrices $\Wl{l} \in \RR^{d^{(l)} \times d^{(l-1)}}$ and bias vectors $b^{(l)} \in \RR^{d^{(l)}}$ for $l = 1, \ldots, L$. On an input $x \in \RR^{d^{(0)}}$ it computes the recursion
\begin{align}
	\actl{0} &\coloneqq x, \notag\\
	\preactl{l} &\coloneqq \Wl{l} \actl{l-1} + b^{(l)}, && l = 1, \ldots, L, \label{eq:forward_recursion}\\
	\actl{l} &\coloneqq \relu\bigl(\preactl{l}\bigr), && l = 1, \ldots, L-1, \notag
\end{align}
where $\relu(z) \coloneqq \max\{z, 0\}$ acts componentwise, and the output layer is affine, $\actl{L} \coloneqq \preactl{L}$. We call $\preactl{l}$ the \emph{pre-activation} and $\actl{l}$ the \emph{(post-)activation} of layer $l$, and denote the \emph{network function} by
\begin{equation}
	f \colon \RR^{d^{(0)}} \longrightarrow \RR^{d^{(L)}},
	\qquad
	f(x) \coloneqq \actl{L}.
	\label{eq:network_function}
\end{equation}
Layers $1, \ldots, L-1$ are the \emph{hidden layers}. Throughout we fix the boundary convention $\relu'(0) \coloneqq 0$, so that the gate indicator $\1\{\preactl{l}_i > 0\}$ returns $0$ at $\preactl{l}_i = 0$.
\end{definition}

For layers $l, \ldots, L$, write $f^{(l)} \colon \RR^{d^{(l)}} \to \RR^{d^{(L)}}$ for the tail map, with $f = f^{(1)} \circ \bigl(x \mapsto \relu(\Wl{1} x + b^{(1)})\bigr)$ and $f^{(L)}$ the affine output map. Appending a constant \emph{intercept input} $1$ absorbs $b^{(l)}$ into a column of $\Wl{l}$; some proofs use this bias-free formulation implicitly.

Chapter~\ref{ch:decompositions} splits parameters and inputs by sign using the following notation.

\begin{definition}[Positive and negative parts]\label{def:posneg_parts}
For a scalar $u \in \RR$ we write $u^{+} \coloneqq \max\{u, 0\}$ and $u^{-} \coloneqq \max\{-u, 0\}$, so that $u = u^{+} - u^{-}$ and $\abs{u} = u^{+} + u^{-}$, with $u^+ u^- = 0$. For a matrix or vector $W$ and a sign label $\sigma \in \{+,-\}$, the \emph{entrywise part} $\posneg{W}{\sigma}$ is obtained by applying $u \mapsto u^{\sigma}$ to every entry. For the layer weights and biases of Definition~\ref{def:ffn_relu} we abbreviate
\begin{equation}
	\Wsig{l}{\sigma} \coloneqq \posneg{\Wl{l}}{\sigma},
	\qquad
	b^{(l,\sigma)} \coloneqq \posneg{b^{(l)}}{\sigma},
	\qquad \sigma \in \{+,-\},
	\label{eq:weight_split}
\end{equation}
so that $\Wl{l} = \Wsig{l}{+} - \Wsig{l}{-}$ and $\abs{\Wl{l}} = \Wsig{l}{+} + \Wsig{l}{-}$ entrywise, and both $\Wsig{l}{\sigma}$ are entrywise non-negative.
\end{definition}

\begin{remark}[Pair notation for decomposed activations]\label{rem:pair_notation}
The activation labels $\act{l}{\sigma}$, $\sigma \in \{+,-\}$, denote the two non-negative \emph{streams} defined in Section~\ref{sec:general_decomposition_def}, distinct from the entrywise parts $\posneg{\actl{l}}{\sigma}$. The streams carry strictly more information than the sign pattern of $\actl{l}$, supplying the sub-threshold signals exploited in Chapter~\ref{ch:saliency}.
\end{remark}

\subsection*{Continuous piecewise linear functions}

The function class computed by \ac{relu} networks admits an exact geometric description.

\begin{definition}[\ac{cpwl} function]\label{def:cpwl}
A function $f \colon \RR^{d} \to \RR^{m}$ is \emph{\acl{cpwl}} (\acs{cpwl}) if it is continuous and there exist finitely many closed polyhedra $Q_1, \ldots, Q_N \subseteq \RR^d$ with $\bigcup_{r=1}^{N} Q_r = \RR^d$ such that the restriction of $f$ to each $Q_r$ is an affine map. The full-dimensional polyhedra on which $f$ is affine, taken maximal with respect to inclusion, are called the \emph{linear regions} of $f$.
\end{definition}

For example, $\abs{\cdot} \colon \RR \to \RR$ is \ac{cpwl} with the two regions $(-\infty, 0]$ and $[0, \infty)$, and $x \mapsto \max\{x_1, \ldots, x_d\}$ is \ac{cpwl} with $d$ regions.

\begin{observation}[\LeanResult{obs:relu_is_cpwl}]\label{obs:relu_is_cpwl}
Every feedforward \ac{relu} network represents a \ac{cpwl} function.\end{observation}

\begin{proof}
Affine maps and componentwise $\relu$ are \ac{cpwl}. Composition preserves this class: if $f$ and $g$ are affine on polyhedral covers $Q_1, \ldots, Q_N$ and $P_1, \ldots, P_M$ of their respective domains, then $g\circ f$ is affine on the cover of at most $NM$ polyhedra $Q_r\cap f^{-1}(P_s)$. These are polyhedra because $f$ is affine on each $Q_r$; composition also preserves continuity. Induction along \eqref{eq:forward_recursion} proves the claim.
\end{proof}

Conversely, every \ac{cpwl} function has an exact \ac{relu} representation with a dimension-dependent depth bound \parencite{Arora2018}. This equivalence underlies the difference-of-convex analysis of Chapter~\ref{ch:decompositions}.

\begin{theoremenv}[\ac{cpwl}--\ac{relu} equivalence, \parencite{Arora2018}, \LeanResult{thm:cpwl_relu_equivalence}]\label{thm:cpwl_relu_equivalence}
A function $f \colon \RR^{d} \to \RR$ is \ac{cpwl} if and only if it is representable by a feedforward \ac{relu} network. Moreover, every \ac{cpwl} function on $\RR^d$ is representable by a \ac{relu} network of depth at most $\bigl\lceil \log_2(d+1) \bigr\rceil + 1$.\end{theoremenv}

\begin{proof}
Observation~\ref{obs:relu_is_cpwl} proves the ``if'' direction. For the converse, every \ac{cpwl} function $f$ admits a \emph{difference-of-convex} representation
\begin{equation}
	f = g - h,
	\qquad
	g = \max_{r \in [N_g]}\, \ell_r,
	\qquad
	h = \max_{s \in [N_h]}\, m_s,
	\label{eq:cpwl_dc_representation}
\end{equation}
with finitely many affine functions $\ell_r, m_s \colon \RR^d \to \RR$ \parencite{Melzer1986_PLasDC,Kripfganz1987_PWAasDC}; their finite maxima $g,h$ are convex. The maximum of two scalars is a two-layer \ac{relu} computation,
\begin{equation}
	\max\{u, v\}
	= \frac{u + v}{2} + \frac{\abs{u - v}}{2}
	= \frac{u + v}{2} + \frac{\relu(u - v) + \relu(v - u)}{2},
	\label{eq:max_as_relu}
\end{equation}
absorbing the affine pre- and post-combinations into surrounding weight matrices. A balanced binary tree of pairwise maxima realizes $g$ at depth $\lceil \log_2 N_g \rceil + 1$, and likewise $h$; a final affine layer forms \eqref{eq:cpwl_dc_representation}, proving representability. For the bound independent of the number of pieces, \textcite{Arora2018} use the Wang--Sun representation of every \ac{cpwl} function on $\RR^d$ as a signed sum of maxima of at most $d+1$ affine functions each. These need at most $\lceil \log_2(d+1) \rceil$ hidden layers and a final affine output layer.
\end{proof}

\begin{remark}[Expressivity and linear regions]\label{rem:linear_regions}
Theorem~\ref{thm:cpwl_relu_equivalence} identifies the function \emph{class}, not the \emph{cost} of representation: deep networks can represent \ac{cpwl} functions with exponentially many linear regions in the number of neurons, an interplay of depth, width, and region counts surveyed by \textcite{huchette2025deeplearningmeetspolyhedral}. The representation \eqref{eq:cpwl_dc_representation} also foreshadows the central theme of this thesis: the convex \ac{cpwl} streams $g$ and $h$ of the DC split can be extracted \emph{constructively and locally} from a trained network, the subject of Section~\ref{sec:dc_splitting_relu}. For a stream $g(x)=\max_r(a_r^\top x+b_r)$, a uniquely active affine piece gives $\nabla g(x)=a_r$: the stream gradient records the slope of its current linear region.
\end{remark}

\section{Convexity and Monotonicity in Neural Networks}\label{sec:convex_monotone_foundations}

Chapter~\ref{ch:decompositions} splits general networks into sign-constrained components. We define the two classes and prove their monotonicity and convexity, respectively.

We order $\RR^d$ coordinatewise: $x \le y$ if and only if $x_i \le y_i$ for all $i \in [d]$. A function $f \colon \RR^d \to \RR^m$ is \emph{monotone} (non-decreasing) if $x \le y$ implies $f(x) \le f(y)$.

\begin{definition}[Monotone neural network]\label{def:monotone_network}
A feedforward network with recursion \eqref{eq:forward_recursion} and monotone activation functions is called a \emph{monotone neural network} if all its weight matrices are entrywise non-negative, $\Wl{l} \ge 0$ for all $l = 1, \ldots, L$.
\end{definition}

Weight-sign constraints are an established device for building monotone models \parencite{Sill1997_MonotonicNetworks,Daniels2010_MonotoneNN} and remain central to modern monotone architectures \parencite{Runje_2023}; monotonicity can also be certified without imposing such constraints \parencite{Liu2020CertifiedMonotone}.

\begin{proposition}[Monotone networks represent monotone functions, \LeanResult{prop:monotone_represents}]\label{prop:monotone_represents}
A monotone neural network represents a monotone function; each coordinate function is non-decreasing in every input variable.\end{proposition}

\begin{proof}
Induct from the monotone input $\actl{0}=x$. Non-negative weights preserve coordinatewise order: $a\le a'$ implies $\Wl{l}a+b^{(l)}\le\Wl{l}a'+b^{(l)}$. Thus monotonicity of $\actl{l-1}$ implies that of $\preactl{l}$; componentwise monotone activation gives monotonicity of $\actl{l}$. The same order argument covers the affine output.
\end{proof}

\begin{definition}[\acl{icnn}, \parencite{Amos2017}]\label{def:icnn}
A feedforward network with recursion \eqref{eq:forward_recursion} whose activation functions are convex \emph{and} monotone is called an \emph{\acl{icnn}} (\acs{icnn}) if all weight matrices except possibly the first are entrywise non-negative, $\Wl{l} \ge 0$ for all $l = 2, \ldots, L$, while $\Wl{1}$ is unconstrained.
\end{definition}

Definition~\ref{def:icnn} covers the plain feedforward core of the architecture of \textcite{Amos2017}; their sign-unconstrained skip (``passthrough'') connections do not affect the statement below, since they enter each pre-activation affinely in $x$.

\begin{proposition}[\acp{icnn} represent convex functions, \LeanResult{prop:icnn_convex}]\label{prop:icnn_convex}
An \ac{icnn} represents a convex function; each coordinate of every activation $\actl{l}$, $l \ge 1$, is a convex function of the input $x$.\end{proposition}

\begin{proof}
We use two closure properties of convex functions: (i) a non-negative linear combination of convex functions plus a constant is convex; (ii) if $\varphi \colon \RR \to \RR$ is convex and non-decreasing and $u \colon \RR^d \to \RR$ is convex, then $\varphi \circ u$ is convex, since for $\lambda \in [0,1]$,
\begin{equation*}
	\varphi\bigl(u(\lambda x + (1-\lambda) y)\bigr)
	\le \varphi\bigl(\lambda u(x) + (1-\lambda) u(y)\bigr)
	\le \lambda\, \varphi(u(x)) + (1-\lambda)\, \varphi(u(y)),
\end{equation*}
where the first inequality uses convexity of $u$ and monotonicity of $\varphi$, and the second convexity of $\varphi$.

Induct over layers. The affine coordinates of $\preactl{1}=\Wl{1}x+b^{(1)}$ are convex, so (ii) gives convexity of $\actl{1}$. If $\actl{l-1}$ has convex coordinates, the non-negative weights for $l\ge2$ make $\preactl{l}=\Wl{l}\actl{l-1}+b^{(l)}$ convex by (i), and (ii) preserves convexity through the activation. The identity activation covers the affine output $\actl{L}=\preactl{L}$.
\end{proof}

\begin{remark}[Monotonicity and convexity combine]\label{rem:monotone_convex_combined}
If, in addition, $\Wl{1} \ge 0$ in Definition~\ref{def:icnn}, the network is simultaneously monotone and input-convex, hence non-decreasing \emph{and} convex by Propositions~\ref{prop:monotone_represents} and \ref{prop:icnn_convex}. Both $\relu$ and Maxout (below) are convex and monotone, so non-negative-weight \ac{relu} networks are of exactly this type. These are the building blocks of the DC splitting of Section~\ref{sec:dc_splitting_relu}: the streams $g$ and $h$ are computed by monotone networks acting on a signed embedding of the input, and are input-convex in the original input. The sign constraint also \emph{limits} expressivity---a weight-constrained \ac{relu} network represents only convex functions, precisely Proposition~\ref{prop:icnn_convex}, which is why monotone architectures add concave counterparts of the activation \parencite{Runje_2023}. The decompositions of this thesis turn the limitation into a feature: the difference of two such networks recovers full generality by Theorem~\ref{thm:cpwl_relu_equivalence} and \eqref{eq:cpwl_dc_representation}.
\end{remark}

\subsection*{Maxout units}

\begin{definition}[Maxout unit and Maxout network, \parencite{Goodfellow2013Maxout}]\label{def:maxout}
A \emph{rank-$k$ Maxout unit} with weights $w_1, \ldots, w_k \in \RR^{d}$ and biases $\beta_1, \ldots, \beta_k \in \RR$ computes
\begin{equation}
	x \longmapsto \maxout(x) \coloneqq \max_{j \in [k]} \bigl( w_j^{\top} x + \beta_j \bigr).
	\label{eq:maxout_unit}
\end{equation}
A \emph{rank-$k$ Maxout network} is a feedforward network whose neurons are rank-$k$ Maxout units; equivalently, layer $l$ carries $k$ weight matrices $W^{(l)}_1, \ldots, W^{(l)}_k$ and biases $b^{(l)}_1, \ldots, b^{(l)}_k$, and
\begin{equation}
	\actl{l}_i \coloneqq \max_{j \in [k]} \bigl( W^{(l)}_j \actl{l-1} + b^{(l)}_j \bigr)_i .
	\label{eq:maxout_layer}
\end{equation}
\end{definition}

\begin{observation}[\ac{relu} and max-pooling as Maxout, \LeanResult{obs:relu_maxout_special_case}]\label{obs:relu_maxout_special_case}
The \ac{relu} activation is the rank-$2$ Maxout unit with parameters $(w_1, \beta_1) = (w, \beta)$ and $(w_2, \beta_2) = (0, 0)$, since $\relu(w^\top x + \beta) = \max\{w^\top x + \beta,\, 0\}$. Max-pooling over a group $J$ of neurons is the rank-$\abs{J}$ Maxout unit whose $j$-th argument is the coordinate projection onto the $j$-th element of $J$. Every rank-$k$ Maxout unit is convex in its input and non-decreasing in each pre-activation argument $w_j^\top x + \beta_j$; hence Maxout qualifies as an admissible activation in Definitions~\ref{def:monotone_network} and \ref{def:icnn}.\end{observation}

\begin{proof}
The first two claims follow from \eqref{eq:maxout_unit}. A maximum of affine functions is convex. For argument-wise monotonicity, set $u_j=w_j^\top x+\beta_j$: $u_j\le u_j'$ for all $j$ implies $\max_j u_j\le\max_j u_j'$.
\end{proof}

\begin{remark}\label{rem:maxout_role}
Maxout networks are the natural piecewise-linear generalization of \ac{relu}: a rank-$k$ unit expresses a maximum of $k$ affine functions in one neuron, matching the building blocks of the DC representation \eqref{eq:cpwl_dc_representation}. The DC splitting of Chapter~\ref{ch:decompositions} is developed for \ac{relu} networks and extended to Maxout units---and to the multiplicative interactions needed for attention---in Chapter~\ref{ch:transformers}, using the exchange identity $\max_{j \in [k]}(u_j - v_j) = \max_{j \in [k]}\bigl\{ u_j + \sum_{i \ne j} v_i \bigr\} - \sum_{i \in [k]} v_i$ for pairs of non-negative streams.
\end{remark}

\section{Game Theory and Markov Decision Processes}\label{sec:game_theory_mdp_foundations}

Chapter~\ref{ch:game_theory} interprets backward propagation as a two-player zero-sum game on the decomposed network graph. We introduce \acp{mdp}, turn-based games, and their optimality equations, restricting throughout to finite state and action spaces.

\begin{definition}[\acl{mdp}]\label{def:mdp}
A (finite) \emph{\acl{mdp}} (\acs{mdp}) \parencite{Puterman1994mdp,SuttonBarto2018rl} is a tuple $\mathcal{M} = (S, A, \P, R, \gamma)$ consisting of a finite state set $S$, a finite action set $A(s)$ for each $s \in S$, a \emph{transition kernel} $\P(s' \mid s, a) \ge 0$ with $\sum_{s'} \P(s' \mid s, a) = 1$, a bounded \emph{reward function} $R(s, a) \in \RR$, and a \emph{discount factor} $\gamma \in [0, 1)$. A \emph{(stationary) policy} is a map $\pi$ assigning to each state $s$ a probability distribution $\pi(\cdot \mid s)$ over $A(s)$; $\pi$ is \emph{pure} (deterministic) if each $\pi(\cdot \mid s)$ is a point mass. A policy induces, for every start state, a Markov chain on $S$ and hence a distribution over trajectories $(s_0, a_0, s_1, a_1, \ldots)$; the \emph{value function} of $\pi$ is the expected discounted return
\begin{equation}
	V^{\pi}(s) \coloneqq \E^{\pi}\Bigl[ \textstyle\sum_{t=0}^{\infty} \gamma^{t}\, R(s_t, a_t) \,\Big|\, s_0 = s \Bigr],
	\label{eq:mdp_value}
\end{equation}
and the \emph{optimal value function} is $V^{\star}(s) \coloneqq \max_{\pi} V^{\pi}(s)$.
\end{definition}

Fixing an enumeration of $S$, a pure policy $\pi$ turns the kernel into a row-stochastic \emph{transition matrix} $P^{\pi} \in \RR^{S \times S}$ with $P^{\pi}_{s s'} = \P(s' \mid s, \pi(s))$ and the reward into a vector $R^{\pi}$, and \eqref{eq:mdp_value} becomes the Neumann series $V^{\pi} = \sum_{t \ge 0} \gamma^t (P^{\pi})^t R^{\pi}$. This matrix picture recurs in Chapter~\ref{ch:game_theory}, where the (sub)stochastic matrices are built from the normalized weight parts $\Wsig{l}{\sigma}$ of Definition~\ref{def:posneg_parts}.

\begin{proposition}[Bellman equations, \LeanResult{prop:bellman}]\label{prop:bellman}
Let $\mathcal{M}$ be a finite \ac{mdp} with $\gamma \in [0,1)$. Then:
\begin{enumerate}
	\item $V^{\pi}$ is the unique solution of the linear system
	\begin{equation}
		V(s) = \sum_{a \in A(s)} \pi(a \mid s) \Bigl[ R(s,a) + \gamma \sum_{s'} \P(s' \mid s, a)\, V(s') \Bigr],
		\qquad s \in S.
		\label{eq:bellman_evaluation}
	\end{equation}
	\item $V^{\star}$ is the unique solution of the \emph{Bellman optimality equation}
	\begin{equation}
		V(s) = \max_{a \in A(s)} \Bigl\{ R(s, a) + \gamma \sum_{s'} \P(s' \mid s, a)\, V(s') \Bigr\},
		\qquad s \in S,
		\label{eq:bellman_optimality}
	\end{equation}
	and any pure policy $\polstar$ that selects, in every state, an action attaining the maximum in \eqref{eq:bellman_optimality} is optimal: $V^{\polstar} = V^{\star}$.
\end{enumerate}\end{proposition}

\begin{proof}
Both claims are standard \parencite{Puterman1994mdp,SuttonBarto2018rl}. The right-hand sides of \eqref{eq:bellman_evaluation} and \eqref{eq:bellman_optimality} define $\gamma$-contractions on $\bigl(\RR^{S}, \norm{\cdot}_{\infty}\bigr)$, so each has a unique fixed point by the Banach fixed-point theorem; conditioning \eqref{eq:mdp_value} on the first transition identifies these fixed points as $V^{\pi}$ and $V^{\star}$, and any policy greedy for $V^{\star}$ is optimal. The acyclic games below make even the contraction argument unnecessary (Remark~\ref{rem:structural_discount}).
\end{proof}

\begin{remark}[State-dependent structural discounts]\label{rem:structural_discount}
Our games use \emph{state-dependent} discounts $\gamma(s)\ge0$ that may exceed $1$. These are \emph{structural} normalizers: for neuron $i$ in layer $l$, the absolute weight mass
\begin{equation}
	\gamma^{(l)}_i \coloneqq \sum_{j} \bigl( \Wsig{l}{+} + \Wsig{l}{-} \bigr)_{ij} = \sum_j \abs{\Wl{l}_{ij}},
	\label{eq:structural_discount}
\end{equation}
restores the mass after normalizing row $i$ of $\abs{\Wl{l}}$ into transition probabilities. Although contraction is unavailable, \emph{layered acyclic} graphs ensure finite trajectories and value sums. Backward induction gives a unique value and optimal policies, possibly nonunique at ties, as the two-player theorem below shows.
\end{remark}

\subsection*{Zero-sum turn-based games}

Two-player games enter because a decomposed network carries two interacting streams: a player $P^{+}$ who collects positive evidence and a player $P^{-}$ who collects negative evidence. We use the following conventions throughout: $p, q, r \in \{+, -\}$ denote player, turn, or sign labels; $\opp{p}$ denotes the opponent ($\opp{+} = -$, $\opp{-} = +$); the \emph{sign-XOR} $p \oplus q$ is $+$ if $p = q$ and $-$ otherwise; and $\xi_{p} \in \{+1, -1\}$ is the numerical sign of the label $p$, so that $\xi_{p}\, \xi_{q} = \xi_{p \oplus q}$.

\begin{definition}[Zero-sum turn-based game]\label{def:turn_based_game}
A (finite) \emph{zero-sum turn-based game} is a tuple $\mathcal{G} = (S_{+}, S_{-}, S_{\mathrm{ter}}, A, \P, R, \gamma)$ where the state set $S = S_{+} \uplus S_{-} \uplus S_{\mathrm{ter}}$ is partitioned into states controlled by player $P^{+}$, states controlled by player $P^{-}$, and terminal states; the action sets $A(s)$, the transition kernel $\P$, and the reward $R$ are as in Definition~\ref{def:mdp} on non-terminal states; $\gamma \colon S \to \RR_{\ge 0}$ is a (possibly state-dependent) discount; and each terminal state $s \in S_{\mathrm{ter}}$ carries a fixed one-time payoff $R(s)$. A \emph{policy pair} $\pi = (\pi_{+}, \pi_{-})$ consists of a policy $\pi_{p}$ for player $P^{p}$ on $S_{p}$. Play generates a trajectory as in an \ac{mdp}, with the player in turn choosing the action; the expected total payoff $V^{\pi}(s)$, defined as in \eqref{eq:mdp_value} with the discount accumulated multiplicatively along the trajectory ($\gamma^t$ replaced by $\prod_{m=0}^{t-1}\gamma(s_m)$), is paid \emph{by $P^{-}$ to $P^{+}$}: player $P^{+}$ maximizes $V^{\pi}$ and player $P^{-}$ minimizes it.
\end{definition}

\begin{theoremenv}[Shapley--Bellman optimality on layered games, \LeanResult{thm:shapley_bellman}]\label{thm:shapley_bellman}
Let $\mathcal{G}$ be a finite zero-sum turn-based game whose transition graph is acyclic. Then $\mathcal{G}$ has a \emph{value}: there is a unique function $V^{\star} \colon S \to \RR$ satisfying the \emph{Shapley--Bellman optimality equations}
\begin{equation}
	V^{\star}(s) =
	\begin{cases}
		\displaystyle \max_{a \in A(s)} \Bigl\{ R(s,a) + \gamma(s) \sum_{s'} \P(s' \mid s, a)\, V^{\star}(s') \Bigr\}, & s \in S_{+}, \\[1.2em]
		\displaystyle \min_{a \in A(s)} \Bigl\{ R(s,a) + \gamma(s) \sum_{s'} \P(s' \mid s, a)\, V^{\star}(s') \Bigr\}, & s \in S_{-}, \\[1.2em]
		R(s), & s \in S_{\mathrm{ter}},
	\end{cases}
	\label{eq:shapley_bellman}
\end{equation}
and there exist pure stationary policies $\polstar = (\polstar_{+}, \polstar_{-})$, obtained by selecting maximizing respectively minimizing actions in \eqref{eq:shapley_bellman}, such that for all $s \in S$ and all policies $\pi_{+}, \pi_{-}$,
\begin{equation}
	V^{(\polstar_{+}, \pi_{-})}(s) \;\ge\; V^{\polstar}(s) = V^{\star}(s) \;\ge\; V^{(\pi_{+}, \polstar_{-})}(s).
	\label{eq:saddle_point}
\end{equation}
In particular, $\polstar$ is a Nash equilibrium of every subgame, i.e.\ a subgame perfect equilibrium.\end{theoremenv}

\begin{proof}
Induct backward along a topological order of the finite acyclic graph. At terminal states, \eqref{eq:shapley_bellman} fixes $V^{\star}(s)=R(s)$ and \eqref{eq:saddle_point} holds. For $s\in S_+$, assume uniqueness and the saddle inequalities at all successors. Each action's bracket in \eqref{eq:shapley_bellman} is finite; maximizing over the finite action set uniquely defines $V^{\star}(s)$ and supplies a pure choice $\polstar_+(s)$. For any opponent policy $\pi_-$, condition on the first transition and apply the successors' left saddle inequality with $\gamma(s)\ge0$:
\begin{align*}
	V^{(\polstar_{+}, \pi_{-})}(s)
	&= R\bigl(s, \polstar_{+}(s)\bigr) + \gamma(s) \sum_{s'} \P\bigl(s' \mid s, \polstar_{+}(s)\bigr)\, V^{(\polstar_{+}, \pi_{-})}(s') \\
	&\ge R\bigl(s, \polstar_{+}(s)\bigr) + \gamma(s) \sum_{s'} \P\bigl(s' \mid s, \polstar_{+}(s)\bigr)\, V^{\star}(s')
	\;=\; V^{\star}(s),
\end{align*}
The last equality uses the maximizing choice $\polstar_+(s)$. Against $\polstar_-$, any $\pi_+$ averages the bracket to at most its maximum; conditioning and the successors' right saddle inequality give $V^{(\pi_+,\polstar_-)}(s)\le V^{\star}(s)$. For $s\in S_-$, exchange $\max$ with $\min$ and reverse the inequalities. Taking $\pi=\polstar$ in both inequalities gives $V^{\polstar}=V^{\star}$. The argument holds in every rooted subgame, proving subgame perfection.
\end{proof}

\begin{remark}\label{rem:shapley_history}
For general discounted stochastic games with cycles, existence of the value and of optimal stationary policies is Shapley's theorem \parencite{Shapley1953stochastic}, proved by the contraction argument of Proposition~\ref{prop:bellman} applied to the \emph{Shapley operator}---the min-max analogue of the Bellman operator \eqref{eq:shapley_bellman}. \textcite{Gaubert2025} formulate their network games in this operator language. The layered acyclic case of Theorem~\ref{thm:shapley_bellman} is all we need, and it additionally delivers subgame perfection, on which Chapter~\ref{ch:trajectory} relies.
\end{remark}

\subsection{The ReLU Net Game of Gaubert and Vlassopoulos}\label{subsec:relu_net_game}

\textcite{Gaubert2025} represent the forward pass of a \ac{relu} network as the value of a zero-sum turn-based stopping game played backward on its graph. We give their full construction in our notation; Chapter~\ref{ch:game_theory} retains its states, transitions, and stopping structure with revised payoff bookkeeping. Fix a network from Definition~\ref{def:ffn_relu} with scalar output ($d^{(L)}=1$), an input $x\in\RR^{d^{(0)}}$, and nonzero weight rows, ensuring strictly positive discounts \eqref{eq:structural_discount}.

Read the identity $\actl{l}_i = \max\bigl\{0,\, \preactl{l}_i\bigr\}$ as a decision: a player at neuron $i$ may \emph{stop}, securing payoff $0$, or \emph{continue}, collecting the bias $b^{(l)}_i$ plus the weighted contributions of the previous layer. Continuing propagates the decision backward through the network; a negative weight flips the sign of its contribution, i.e.\ hands the turn to an adversary facing the mirrored problem. The optimal value at neuron $i$ is exactly the activation $\actl{l}_i$.

\begin{definition}[ReLU Net Game, following \parencite{Gaubert2025}]\label{def:relu_net_game}
The \emph{ReLU Net Game} on the network $f$ at input $x$ is the zero-sum turn-based game with the following components.
\begin{itemize}
	\item \textbf{States.} For every hidden layer $l \in \{1, \ldots, L-1\}$, every neuron $i \in [d^{(l)}]$, and every turn label $q \in \{+,-\}$ there is an \emph{activation state} $s^{(l)}_{i,q}$, controlled by player $P^{q}$, so that $S_q = \bigl\{ s^{(l)}_{i,q} : l \in [L-1],\, i \in [d^{(l)}] \bigr\}$. For every input coordinate $k \in [d^{(0)}]$ and every $q \in \{+,-\}$ there is a \emph{terminal state} $s^{(0)}_{k,q}$, and there is one absorbing \emph{cemetery state} $\perp$ with zero payoff. Finally, there is a single \emph{output state} $s^{(L)}$ at the affine output neuron, at which no decision is taken.
	\item \textbf{Actions.} At every activation state the player in turn chooses an action from $\{\mathrm{stop}, \mathrm{cont}\}$. The output state has the single action $\mathrm{cont}$.
	\item \textbf{Transitions.} Stopping moves deterministically to the cemetery, $\P(\perp \mid s^{(l)}_{i,q}, \mathrm{stop}) = 1$. Continuation at $s^{(l)}_{i,q}$ moves one layer backward, to neuron $j$ of layer $l-1$ with probability proportional to the absolute weight, keeping the turn on positive weights and passing it to the opponent on negative weights:
	\begin{equation}
		\P\bigl( s^{(l-1)}_{j,r} \,\big|\, s^{(l)}_{i,q}, \mathrm{cont} \bigr)
		\coloneqq \frac{\Wsig{l}{q \oplus r}_{ij}}{\gamma^{(l)}_i},
		\qquad
		\gamma^{(l)}_i = \sum_{j} \bigl( \Wsig{l}{+} + \Wsig{l}{-} \bigr)_{ij},
		\label{eq:rng_kernel}
	\end{equation}
	where for $l = 1$ the successor $s^{(0)}_{j,r}$ is the corresponding terminal state. The output state transitions by the same rule \eqref{eq:rng_kernel} with $l = L$ and turn label $q = +$.
	\item \textbf{Payoffs and discounts.} Continuation at $s^{(l)}_{i,q}$ pays the signed bias $\xi_{q}\, b^{(l)}_i$; stopping pays $0$. A terminal state pays the signed input coordinate once, $R\bigl(s^{(0)}_{k,q}\bigr) \coloneqq \xi_{q}\, x_k$, and the game ends. The output state pays $b^{(L)}$ upon its forced continuation. The discount at every activation state and at the output state is the structural discount $\gamma\bigl(s^{(l)}_{i,q}\bigr) \coloneqq \gamma^{(l)}_i$ of \eqref{eq:structural_discount}.
\end{itemize}
Player $P^{+}$ maximizes and $P^{-}$ minimizes the expected discounted total payoff, as in Definition~\ref{def:turn_based_game}.
\end{definition}

The game graph is layered and acyclic---every transition strictly decreases the layer index or ends in the cemetery---so Theorem~\ref{thm:shapley_bellman} applies despite the discounts $\gamma^{(l)}_i$ being arbitrary positive numbers. Normalization and discounting cancel: the product of the transition probability \eqref{eq:rng_kernel} and the discount $\gamma^{(l)}_i$ is exactly the absolute weight $\abs{\Wl{l}_{ij}}$, so the Shapley--Bellman recursion reproduces the affine algebra of the forward pass.

\begin{theoremenv}[Forward pass as game value, \parencite{Gaubert2025}, \LeanResult{thm:relu_net_game_value}]\label{thm:relu_net_game_value}
The ReLU Net Game of Definition~\ref{def:relu_net_game} admits the subgame perfect equilibrium $\polstar = (\polstar_{+}, \polstar_{-})$ in which both players use the same pure stopping rule,
\begin{equation}
	\polstar_{q}\bigl( s^{(l)}_{i,q} \bigr) = \mathrm{cont}
	\quad \Longleftrightarrow \quad
	\preactl{l}_i > 0,
	\label{eq:rng_equilibrium_policy}
\end{equation}
and the value function of the game evaluates the network:
\begin{equation}
	V^{\star}\bigl(s^{(l)}_{i,q}\bigr) = \xi_{q}\, \actl{l}_i
	\quad \text{for } l \in \{0, 1, \ldots, L-1\},
	\qquad
	V^{\star}\bigl(s^{(L)}\bigr) = f(x),
	\label{eq:rng_value_identity}
\end{equation}
where for $l = 0$ we read $\actl{0}_k = x_k$. The equilibrium is unique up to tie-breaking at neurons with $\preactl{l}_i = 0$, where both actions are optimal; the convention \eqref{eq:rng_equilibrium_policy} stops at ties, matching $\relu'(0) = 0$.\end{theoremenv}

\begin{proof}
By Theorem~\ref{thm:shapley_bellman}, it suffices to check \eqref{eq:rng_value_identity} and policy \eqref{eq:rng_equilibrium_policy} against \eqref{eq:shapley_bellman}. Induct on layers.

\emph{Base case $l=0$.} The terminal payoff $R\bigl(s^{(0)}_{k,q}\bigr)=\xi_q x_k=\xi_q\actl{0}_k$ satisfies \eqref{eq:rng_value_identity}.

\emph{Inductive step.} Assume \eqref{eq:rng_value_identity} at layer $l-1$, $l\ge1$. At $s^{(l)}_{i,+}\in S_+$, stopping yields $0+V^{\star}(\perp)=0$; by \eqref{eq:rng_kernel} and induction, continuation yields
\begin{align}
	R\bigl(s^{(l)}_{i,+}, \mathrm{cont}\bigr) + \gamma^{(l)}_i \sum_{j, r} \frac{\Wsig{l}{+ \oplus r}_{ij}}{\gamma^{(l)}_i}\, \xi_{r}\, \actl{l-1}_j
	&= b^{(l)}_i + \sum_{j} \Bigl( \Wsig{l}{+}_{ij}\, \actl{l-1}_j - \Wsig{l}{-}_{ij}\, \actl{l-1}_j \Bigr) \notag\\
	&= b^{(l)}_i + \sum_j \Wl{l}_{ij}\, \actl{l-1}_j
	\;=\; \preactl{l}_i,
	\label{eq:rng_cont_value}
\end{align}
using $+ \oplus r = r$, $\xi_{+} = 1$, $\xi_{-} = -1$, and $\Wl{l} = \Wsig{l}{+} - \Wsig{l}{-}$. The maximizing player $P^{+}$ therefore obtains
\begin{equation*}
	V^{\star}\bigl(s^{(l)}_{i,+}\bigr) = \max\bigl\{ 0,\, \preactl{l}_i \bigr\} = \relu\bigl(\preactl{l}_i\bigr) = \actl{l}_i = \xi_{+}\, \actl{l}_i,
\end{equation*}
continuing exactly when $\preactl{l}_i > 0$. At the mirrored state $s^{(l)}_{i,-} \in S_{-}$ the continuation payoff is $\xi_{-} b^{(l)}_i = -b^{(l)}_i$, and the kernel labels satisfy $- \oplus r = \opp{r}$, so the computation \eqref{eq:rng_cont_value} returns the continuation value
\begin{equation*}
	-b^{(l)}_i + \sum_{j} \Bigl( \Wsig{l}{-}_{ij}\, \actl{l-1}_j - \Wsig{l}{+}_{ij}\, \actl{l-1}_j \Bigr)
	= -\preactl{l}_i .
\end{equation*}
Thus $P^-$ obtains $V^{\star}\bigl(s^{(l)}_{i,-}\bigr)=\min\{0,-\preactl{l}_i\}=-\relu\bigl(\preactl{l}_i\bigr)=\xi_-\actl{l}_i$, preferring continuation exactly when $-\preactl{l}_i<0$, the gate \eqref{eq:rng_equilibrium_policy}. At $\preactl{l}_i=0$, both actions give either player value $0$, proving uniqueness up to ties. Forced continuation at the affine output gives, by \eqref{eq:rng_cont_value} with $l=L$, $V^{\star}\bigl(s^{(L)}\bigr)=\preactl{L}=\actl{L}=f(x)$.
\end{proof}

\begin{remark}[Softplus networks from entropy regularization]\label{rem:softplus_game}
\textcite{Gaubert2025} also treat softplus networks, by \emph{regularizing} the stop/continue decision. Replace the hard choice at a positive activation state by a mixed policy $\nu \in \Delta(\{\mathrm{stop}, \mathrm{cont}\})$ rewarded with a Shannon-entropy bonus at temperature $\theta > 0$. At a negative state, the entropy term enters with the opposite sign in the payoff to $P^+$. By the Legendre--Fenchel identity for the entropy \parencite{BoydVandenberghe2004},
\begin{equation}
	\max_{\nu \in \Delta_n} \Bigl\{ \textstyle\sum_{i=1}^{n} \nu_i r_i + \theta\, \mathcal{H}(\nu) \Bigr\}
	= \theta \log \textstyle\sum_{i=1}^{n} e^{r_i / \theta},
	\qquad
	\mathcal{H}(\nu) = -\textstyle\sum_i \nu_i \log \nu_i,
	\label{eq:legendre_fenchel}
\end{equation}
with the Gibbs distribution $\nu^{\star}_i \propto e^{r_i/\theta}$ as unique maximizer, the regularized value at a positive state with continuation value $\preactl{l}_i$ (computed as in \eqref{eq:rng_cont_value}) becomes $\theta \log\bigl(1 + e^{\preactl{l}_i / \theta}\bigr)$---the softplus activation at temperature $\theta$---and the equilibrium policy continues with the sigmoid probability $\bigl(1 + e^{-\preactl{l}_i/\theta}\bigr)^{-1}$. The hard game of Definition~\ref{def:relu_net_game} is recovered in the limit $\theta \to 0$. Chapter~\ref{ch:game_theory} re-derives this mechanism from the regularized decompositions of Section~\ref{sec:regularized_decompositions} and reuses it at the routing states of the \ac{rg}, with the two temperatures $\theta$ (activation gates) and $\tau$ (routing) kept distinct.
\end{remark}

\begin{remark}[What Chapter~\ref{ch:game_theory} inherits and what it changes]\label{rem:inherited_vs_new}
The \ac{sg} of Chapter~\ref{ch:game_theory} keeps the structural skeleton of Definition~\ref{def:relu_net_game}: per-neuron $\pm$ states, stop-versus-continue actions, the bias as continuation payoff and the input as final payoff, the transition kernel \eqref{eq:rng_kernel} driven by the weight parts $\Wsig{l}{\sigma}$, the cemetery state, and the structural discounts $\gamma^{(l)}_i$. Two changes are load-bearing for attribution. First, the original tracks only the single signed value \eqref{eq:rng_value_identity}; Chapter~\ref{ch:game_theory} additionally maintains a \emph{non-negative payoff per player}, identifying the pair of player-specific values $\Vp{p}{s}$, $p \in \{+,-\}$, with the two streams of the Stopping Decomposition of Section~\ref{sec:stopping_decomposition} and recovering \eqref{eq:rng_value_identity} as their difference. Second, the signed terminal payoffs $\xi_q\, x_k$ become the non-negative pair $(x_k^{+}, x_k^{-})$ of Definition~\ref{def:posneg_parts}, keeping every player-specific quantity non-negative and exposing the parity structure of backward trajectories. Chapter~\ref{ch:game_theory} further extends the graph to skip connections and max-pooling and adds risk aversion and regularized policies, none of which appear in the original game.
\end{remark}
